\documentclass[11pt,a4paper]{article}
\usepackage[utf8]{inputenc}
\usepackage[T1]{fontenc}
\usepackage[margin=0.85in]{geometry}
\usepackage{amsmath,amssymb}
\usepackage{booktabs}
\usepackage{multirow}
\usepackage{longtable}
\usepackage{array}
\usepackage{enumitem}
\usepackage{xcolor}
\usepackage[hidelinks]{hyperref}
\usepackage{titlesec}
\usepackage{parskip}

\definecolor{qnavy}{RGB}{20,40,75}
\definecolor{qgrey}{RGB}{90,90,90}
\definecolor{qgreen}{RGB}{20,110,70}
\definecolor{qred}{RGB}{150,30,30}
\titleformat{\section}{\large\bfseries\color{qnavy}}{\thesection}{0.6em}{}
\titleformat{\subsection}{\normalsize\bfseries\color{qnavy}}{\thesubsection}{0.6em}{}

\newcommand{\ver}{\textcolor{qgreen}{\textbf{V}}}
\newcommand{\src}{\textcolor{qgrey}{\textbf{S}}}
\newcommand{\gap}{\textcolor{qred}{\textbf{--}}}

\title{\vspace{-1.6em}\color{qnavy}\bfseries Data Inventory\\[4pt]
\large\normalfont\color{qgrey}Quantara three-study data programme --- glioma, hepatocellular
carcinoma, lung}
\author{\normalsize Quantara \;\textperiodcentered\; Reza Nehzati \;\textperiodcentered\;
\texttt{reza@heydonto.com}}
\date{\normalsize 4 August 2026 \;\textperiodcentered\; Internal --- team reference}

\begin{document}
\maketitle
\thispagestyle{empty}

\begin{abstract}
\noindent
This document inventories the corpus assembled for the three-study programme and held in
Google Cloud Storage. It totals approximately \textbf{2.16\,TiB}: a DICOM imaging layer of 2{,}031 patients
across 2.13 million objects, five spatial/multiplex imaging cohorts, eighteen molecular and
clinical label datasets, and --- added since the 4 August edition --- a whole-slide-imaging
layer of 1{,}053 TCGA lung slides with their encoded tile features and an assembled
TCGA lung methylation matrix.
Every item was screened for licence before acquisition and no item under a
non-commercial restriction was taken. Counts marked \ver{} were verified by comparing what
arrived against what was planned; counts marked \src{} come from the source publication
and were not independently confirmed. That distinction is deliberate and is explained in
\S2.

\medskip
\noindent\textbf{Revision, 16 August 2026.} Layer~D gains an external whole-slide
holding (\S5) --- 133 slides from two institutions outside TCGA, acquired to test whether the
site signature reported in R12--R15 reproduces across archives --- and Layer~C gains the
ACRIN 6668 clinical archive used in R16. One row of Layer~B is \textbf{corrected}: the LUSC
imaging mass-spectrometry deposit holds four images, not 330 patients (\S4). Sizes for the new
items were measured directly on 16 August; earlier layers were not re-measured.

\medskip
\noindent\textbf{Revision, 14 August 2026.} Layers A--C are unchanged from the 4 August
edition and their object counts are \emph{carried forward} from that verification rather than
re-walked --- the DICOM prefix holds over two million objects and a fresh enumeration costs
more than it adds when nothing was written to it. What is new is Layer D (\S5), measured
directly on 14 August. Readers who need Layer A re-verified should say so; the exact-match
check is reproducible from \texttt{\_provenance/}.
\end{abstract}

\section{Location and access}

\begin{itemize}[leftmargin=1.4em,nosep]
\item \textbf{Bucket:} \texttt{gs://heydonto-quantara-lungcdx/data-request-2026-08/}
(project \texttt{heydonto-425716}, region \texttt{US-CENTRAL1}).
\item \textbf{Layout:} \texttt{imaging/} (DICOM, by disease), \texttt{glioma/spatial/},
\texttt{lung/mif/}, \texttt{lung/msi/}, \texttt{hcc/imaging/}, \texttt{labels/} (by
disease), \texttt{\_provenance/}, \texttt{\_run/}.
\item \textbf{Provenance:} \texttt{\_provenance/} holds the acquisition manifest with a
per-item licence and the exclusion list; \texttt{\_run/} holds the acquisition job, the
transfer log and per-collection completion markers. The pull is re-runnable and
idempotent.
\item \textbf{Read-only external access} can be granted per-prefix; see \S8.
\end{itemize}

\section{How to read this document}

Two symbols appear throughout, because during assembly several transfers reported success
while delivering little or nothing, and several published patient counts did not survive
checking.

\begin{itemize}[leftmargin=1.4em,nosep]
\item \ver{} --- \textbf{verified by us.} Delivered file or object counts were compared
against planned counts for that dataset. Where a source of truth existed (the imaging
index, a cohort's own metadata table), the values were recomputed from it.
\item \src{} --- \textbf{from the source publication.} Typically patient or subject
totals, where the deposited metadata carries no such field. The files are verified; the
subject count is taken on trust.
\end{itemize}

\noindent
Sizes are binary (GiB/TiB) as reported by object storage. Patient counts and series
counts for the DICOM layer were recomputed from the imaging index rather than taken from
collection pages, which are known to overstate what actually shipped.

\section{Layer A --- DICOM imaging}

Copied server-side from the public imaging commons. All ten collections complete:
\textbf{2{,}091{,}646 of 2{,}091{,}646 expected DICOM instances}, an exact match with no
series or instance dropped. \ver

\begin{center}
\small
\begin{longtable}{@{}lrrrrl@{}}
\toprule
\textbf{Collection} & \textbf{Patients} & \textbf{Series} & \textbf{Objects} &
\textbf{GiB} & \textbf{Licence} \\
\midrule
\endfirsthead
\toprule \textbf{Collection} & \textbf{Patients} & \textbf{Series} & \textbf{Objects} &
\textbf{GiB} & \textbf{Licence} \\ \midrule
\endhead
\bottomrule
\endfoot
\multicolumn{6}{@{}l}{\itshape Glioma} \\
UPENN-GBM (MR $+$ segmentations) & 630 & 6{,}064 & 830{,}618 & 138.0 & CC BY 4.0 \\
CPTAC-GBM (pathology WSI only) & 178 & 462 & 2{,}835 & 98.6 & CC BY 4.0 \\
\addlinespace
\multicolumn{6}{@{}l}{\itshape Lung} \\
NSCLC-Radiogenomics (CT/PET $+$ EGFR/KRAS) & 211 & 1{,}722 & 287{,}125 & 97.4 & CC BY 3.0 \\
Lung-PET-CT-Dx & 355 & 2{,}606 & 252{,}446 & 124.6 & CC BY 4.0 \\
ACRIN-NSCLC-FDG-PET & 242 & 3{,}762 & 497{,}752 & 142.2 & CC BY 3.0/4.0 \\
Anti-PD-1-Lung (2 timepoints) & 46 & 819 & 134{,}607 & 60.3 & CC BY 3.0 \\
LungCT-Diagnosis & 61 & 61 & 4{,}682 & 2.4 & CC BY 3.0 \\
Lung-Fused-CT-Pathology & 6 & 52 & 11{,}210 & 5.8 & CC BY 3.0 \\
\addlinespace
\multicolumn{6}{@{}l}{\itshape Liver} \\
HCC-TACE-Seg (CT $+$ segmentations) & 105 & 1{,}020 & 52{,}311 & 29.1 & CC BY 4.0 \\
Colorectal-Liver-Metastases & 197 & 618 & 18{,}060 & 11.4 & CC BY 4.0 \\
\midrule
\textbf{Total} & \textbf{2{,}031} & \textbf{17{,}186} & \textbf{2{,}091{,}646} &
\textbf{709.8} & \\
\end{longtable}
\end{center}

\noindent
All 26{,}258 candidate series were licence-screened before acquisition; none carried a
non-commercial restriction. \ver{} Two properties worth noting for study design:
\textbf{CPTAC-GBM contains no MRI} (its 462 series are all slide microscopy), and
\textbf{UPENN-GBM contains no clinical table} in the DICOM layer --- its molecular
annotations are a separate file, held in Layer~C.

\section{Layer B --- spatial, multiplex and non-DICOM imaging}

\begin{center}
\small
\begin{longtable}{@{}p{5.1cm}rrp{3.0cm}p{3.4cm}@{}}
\toprule
\textbf{Dataset} & \textbf{Files} & \textbf{GiB} & \textbf{Licence} & \textbf{Content} \\
\midrule
\endfirsthead
\toprule \textbf{Dataset} & \textbf{Files} & \textbf{GiB} & \textbf{Licence} &
\textbf{Content} \\ \midrule
\endhead
\bottomrule
\endfoot
Glioma spatial multi-omics & 4{,}599 \ver & 335.6 & \textbf{CC0} & Multiplexed ion-beam
imaging $+$ spatial transcriptomics $+$ glycomics; 310 patients \src \\
NSCLC multiplex immunofluorescence & 1{,}097 \ver & 103.8 & EBI default (no additional
restriction) & PD-L1 / CD3 / CD8 antibody comparison; 431 patients \src \\
GBM mpMRI, multi-centre & 7{,}071 \ver & 41.5 & \textbf{CC0} & T1/T1CE/T2/FLAIR;
\textbf{337 participants} \ver, MGMT on all 337 \ver, \textbf{21 distinct scanner
models} \ver \\
HCC multiphase CT (WAW-TACE) & 10 \ver & 44.9 & CC BY 4.0 & Native/arterial/portal/delayed
CT $+$ 377 tumour masks $+$ radiomics; 233 treatment-naive patients \src \\
LUSC imaging mass spectrometry & 7 \ver & 14.5 & CC BY 4.0 & \textcolor{qred}{\textbf{Correction
2026-08-16:}} the deposit (S-BIAD814) is \textbf{four} \texttt{.czi} whole-slide
MALDI/cytokeratin--vimentin overlay images plus three manifest files --- \emph{not}
per-patient data for 330 patients. Its own \texttt{Files\_list.tsv} lists exactly four
images \ver. The 330-patient figure was the publication's cohort, not the deposit's
contents; the same error class as the anti-PD-1 methylation row (\S5.3). Usable as imaging,
not as a cohort \\
\midrule
\textbf{Total} & \textbf{12{,}784} & \textbf{540.3} & & \\
\end{longtable}
\end{center}

\section{Layer C --- molecular and clinical labels}

Eighteen datasets, approximately 27\,GiB. Marker counts below were recomputed from the
delivered files. \ver

\subsection*{5.1 Glioma}

\begin{center}
\small
\begin{tabular}{@{}p{4.0cm}p{10.4cm}@{}}
\toprule
\textbf{Dataset} & \textbf{Content} \\
\midrule
MSK glioma panel & 923 patients / 1{,}004 samples. TERT promoter status determinable for
all 1{,}004 (every panel used includes TERT), with \textbf{596 promoter-mutation records
across 594 samples}; \textbf{302 CDKN2A homozygous deletions}. Plus WHO grade, prior
therapy lines, TMB, treatment timeline. \emph{No imaging.} \\
TCGA pan-glioma marker table & $n=1{,}122$. TERT promoter status on \textbf{329}
(162 mutant / 167 wild-type); MGMT promoter 932; IDH 995; IDH/1p19q subtype 1{,}093
(171 co-deleted); methylation and expression clusters. \emph{No imaging.} \\
UPENN-GBM clinical & The \textbf{same 630 patients as the imaging layer} (exact join
after normalising a timepoint suffix). IDH1 recorded for all 630, of which \textbf{533
definite} (wild-type/mutated) and 97 NOS/NEC; \textbf{MGMT definite on 275}; survival on
630. \\
UTSW-Glioma metadata & 625 patients, adult diffuse glioma grades 2--4. IDH \textbf{622}
(\emph{413 by immunohistochemistry only}, which detects IDH1-R132H alone); 1p/19q
\textbf{335}; MGMT \textbf{281}. \emph{Labels only --- the images moved behind controlled
access.} \\
\bottomrule
\end{tabular}
\end{center}

\subsection*{5.2 Hepatocellular carcinoma}

\begin{center}
\small
\begin{tabular}{@{}p{4.0cm}p{10.4cm}@{}}
\toprule
\textbf{Dataset} & \textbf{Content} \\
\midrule
Chinese Liver Cancer Atlas & 494 samples; RNA-seq TPM on 238, with FASN present in both
raw and $z$-scored matrices. Independently reproduced here: FASN co-expression
$\rho=+0.546$ ACACA, $+0.510$ SCD, $+0.495$ SREBF1. \ver \\
HCC lncRNA arrays & 400 arrays (200 tumour, 200 matched non-cancerous),
prognosis-annotated. \\
\bottomrule
\end{tabular}
\end{center}

\noindent
The TCGA-LIHC layer (371 primary tumours, 50 solid-tissue normals with 50 matched pairs;
FASN protein on 184 by RPPA; 165 tumour/normal proteomic pairs) is reachable openly and
was used for analysis, but is not mirrored into this bucket.

\subsection*{5.3 Lung}

\begin{center}
\small
\begin{tabular}{@{}p{4.0cm}p{10.4cm}@{}}
\toprule
\textbf{Dataset} & \textbf{Content} \\
\midrule
MSK NSCLC ctDx & 2{,}621 patients; circulating-tumour-DNA panel $+$ treatment timeline. \\
MSK LUAD & 604 patients with treatment context. \\
Post-osimertinib-failure cohort & \textbf{183 patients}, EGFR-mutant NSCLC after
first-line osimertinib failure; panel genomics and clinical course. \emph{See licence
caveat, \S7.} \\
Cell-line methylation & \textbf{9.02\,GB} archive (GSE68379), GEO declares 1{,}028 lines on
450k linked by COSMIC identifier to dose--response. \textcolor{qred}{\textbf{Correction
2026-08-20:} our mirror is \emph{silently corrupt}. It is byte-complete against both GCS and GEO's
declared size (9{,}020{,}579{,}840), ends with a valid tar terminator and raises no error, yet
enumerates \textbf{1{,}302 of 2{,}056 members} --- 649 of 1{,}028 cell lines, 127 of 198 lung. Two
extraction modes gave the identical truncation. R23 bypassed it and fetched the 198 lung samples
per-sample from GEO (396/396, zero failures). The 1{,}028 figure describes GEO's copy, not ours.
\textbf{Every other bulk archive in this inventory should be assumed suspect until its members are
counted}, since size and checksum both passed here. A processed matrix now exists at
\texttt{.../GSE68379/processed/}: \textbf{485{,}512 probes $\times$ 198 lung cell lines}
(313\,MB \texttt{lung\_betas.npz} plus \texttt{lung\_meta.csv} with cell line, histology and COSMIC
identifier), built from IDATs fetched per sample rather than from the archive. 187 of those lines
carry GDSC2 dose--response; report R23 uses them.} \\
EGFR-TKI-treated methylation & 79 advanced LUAD, all TKI-treated, 450k $+$
objective-response and disease-control endpoints. \\
Anti-PD-1 methylation & 81 patients, EPIC array. \textcolor{qred}{\textbf{Correction
2026-08-10:}} the deposit (GSE115246, EPIMMUNE) carries \emph{no outcome labels} ---
metadata is disease status / tissue / cohort only; responder status exists solely in the
paywalled journal appendix. Arrays verified real (signal supplement with detection
$p$-values present; no IDATs, processed betas only). Previously described here as
``responder / non-responder'' --- that was the publication's design, not the deposit's
contents. \\
Paired pre/post-osimertinib & \textbf{4 patients} sampled before and after osimertinib,
profiled by \emph{spatial} transcriptomics (Visium~HD with StarDist segmentation). To our
knowledge the only public paired pre/post-progression tissue resource in EGFR-mutant
NSCLC. \\
SWI/SNF resistance series & PC-9 and its osimertinib-resistant derivative:
\textbf{11 of 14} accessibility tracks (three failed transfer, \S6); matched RNA
expression; H1975 / HCC827 copy-number and variant calls. \\
Persister chromatin & PC-9 (osimertinib) and H358 (sotorasib) persisters,
single-cell ATAC. \\
ACRIN 6668 clinical & \textbf{Added 2026-08-16.} 251 registered patients, stage-III NSCLC
treated with chemoradiotherapy; core-lab PET assessments, survival, staging, performance
status and follow-up across 15 case-report forms (two release samples, 525\,kB). Distributed
by TCIA directly under CC BY 3.0 --- \emph{no} NCTN application. This is the cohort of
report R16. \\
\bottomrule
\end{tabular}
\end{center}

\subsection*{5.4 Derived cohort available from Layer C}

Verified by intersection rather than asserted: \textbf{137 KMT2D, 52 SMARCA4 and 133
KEAP1} chromatin-regulator-mutant patients with methylome, diagnostic-slide morphology
\emph{and} survival all present, plus \textbf{125} patients with matched tumour/normal
methylation. These sit correctly below independently derived mutant totals of 152, 63 and
151 respectively. \ver

\section{Layer D --- whole-slide imaging and derived features}

Added between 10 and 14 August 2026 to support the NSCLC whole-slide work (reports R12--R15).
All three items were measured directly on 14 August.

\begin{center}
\begin{tabular}{@{}llrl@{}}
\toprule
\textbf{Item} & \textbf{Contents} & \textbf{Size} & \textbf{Licence} \\
\midrule
\texttt{wsi/TCGA-LUAD} & 541 diagnostic slides & \multirow{2}{*}{767.1 GiB} & TCGA open \ver{} \\
\texttt{wsi/TCGA-LUSC} & 512 diagnostic slides & & TCGA open \ver{} \\
\texttt{features/phikon-v2/} & tile embeddings: TCGA, CPTAC, Vanguri PD-L1 & 59.9 GiB & derived \ver{} \\
\texttt{lung/tcga\_methylation\_450k/} & assembled 450k beta matrix & 1.7 GiB & TCGA open \ver{} \\
\midrule
\multicolumn{4}{@{}l}{\itshape Added 16 August 2026 --- external institutions, outside TCGA} \\
\texttt{wsi\_external/} & 133 slides: EAGLE (49) and HTMCP-LC (84) & 101.1 GiB & open \ver{} \\
\texttt{features/phikon-v2/external/} & tile embeddings for 129 of those slides & 1.6 GiB & derived \ver{} \\
\bottomrule
\end{tabular}
\end{center}

Four things about this layer are worth stating because they are not obvious from the sizes.

\textbf{The slides were processed, not merely stored.} The 1{,}053 TCGA slides were tiled into
14{,}721{,}903 tiles at approximately 0.5\,\textmu m/px with an Otsu tissue filter, then encoded
with Phikon-v2. The per-slide manifest at
\texttt{features/phikon-v2/tcga/\_manifest.csv} records tile count, measured microns-per-pixel
and whether that value was read from the file or assumed, for every slide.

\textbf{CPTAC slides are not retained --- only their features.} The CPTAC-LUAD/LSCC set
(432 patients, 1{,}343 slides) was tiled and encoded, and the encoded features are held, but the
source slides were not kept. They are re-obtainable from NCI on demand, and storing a second
copy of a public archive was not worth 500-odd GiB.

\textbf{The methylation matrix is derived, and deliberately so.} It represents 919 GDC
files totalling 11.23 GiB of source data, assembled into a single
486{,}427\,$\times$\,919 float32 matrix. The raw files were parsed and discarded; GDC is a
permanent archive keyed by \texttt{file\_id}, so \texttt{MANIFEST\_450k.tsv} with its MD5s is
sufficient provenance to reconstitute them. This is why 11.23 GiB of input appears as 1.7 GiB of
stored object, which would otherwise look like a transfer failure.

\textbf{Four of the 133 external slides did not encode, and the reason matters more than the
count.} The four failures returned \texttt{empty\_no\_tissue} under the tissue filter: they are
immunohistochemistry, not H\&E. That is a property of HTMCP-LC generally --- only 30 of its 84
slides are H\&E. Any analysis that pools all 84 and calls the result an \emph{institutional}
signature is at risk of measuring a \emph{stain} signature instead, so the H\&E subset is the
unit to use and the manifest records stain per slide.

\section{Known gaps}

\subsection*{6.1 Residual transfer gaps on our side}
\begin{itemize}[leftmargin=1.4em,nosep]
\item Three of fourteen accessibility tracks in the SWI/SNF series (a size-check bug on
our side, not an availability problem).
\item One 198-slide lung cohort with per-slide molecular group labels is catalogued but
not mirrored at any public path.
\end{itemize}

\subsection*{6.2 Available only behind a permission gate}
Obtainable by asking, not by downloading:
\begin{itemize}[leftmargin=1.4em,nosep]
\item A 192-patient Central European cohort with eight-sequence MRI \emph{plus} TERT
promoter and CDKN2A --- the only dataset located that pairs multiparametric MRI with both
markers. Non-commercial licence; permits commercial use by explicit written permission.
It has no MGMT, and only 156 of its 192 examinations are pre-treatment.
\item A pan-cancer registry ($\sim$228k patients) and an immunotherapy cohort with PD-L1
immunohistochemistry (247 patients) --- both click-through registrations not yet completed.
\item Glioma MRI at UCSF, UT~Southwestern, Missouri and Río~Hortega, which moved behind
controlled access during 2025. We hold UT~Southwestern's molecular labels but not its
images.
\end{itemize}

\subsection*{6.3 Not obtainable at any price}
\begin{itemize}[leftmargin=1.4em,nosep]
\item A resolvable genomic locus for the transcript designated HELC.
\item FASN enzymatic activity in human HCC tissue at any scale approaching 100 patients;
it is not recoverable from archival material.
\item A glioma cohort with all five markers \emph{and} multiparametric MRI beyond
\textbf{eight} patients. Four MRI sequences plus IDH reaches roughly 1{,}840 patients and
adding 1p/19q about 944, but the five-marker intersection collapses to eight.
\end{itemize}

\section{Licence register}

\begin{center}
\small
\begin{tabular}{@{}p{3.6cm}p{4.4cm}p{6.4cm}@{}}
\toprule
\textbf{Licence} & \textbf{Applies to} & \textbf{Commercial position} \\
\midrule
CC0 & Glioma spatial multi-omics; multi-centre GBM mpMRI & Unrestricted \\
CC BY 3.0 / 4.0 & All ten DICOM collections; HCC multiphase CT; LUSC mass spectrometry;
UTSW and UPENN metadata & Permitted, with attribution \\
No additional restriction & NSCLC multiplex immunofluorescence & Permitted (repository
default) \\
Open Database Licence & MSK glioma panel; MSK NSCLC ctDx; Chinese Liver Cancer Atlas &
Commercial use \emph{explicitly} permitted. Share-alike attaches only to
\textbf{derivative databases}: a trained model is a ``produced work'' and is not itself
copyleft, but a curated training table is a derivative database and must be offered on
the same terms. Legal read required before any data product ships. \\
Broad GDAC / NIH terms & TCGA pan-glioma marker table & Permitted; no restriction on
commercial product development \\
NCBI (no restrictions) & All expression, methylation and chromatin series & Permitted;
residual submitter-IP caveat \\
\textbf{No licence stated} & Post-osimertinib-failure cohort & Open access, direct
download, no data-use agreement --- but commercial use is addressed in neither direction.
\textbf{No affirmative grant; do not build product on it, and do not redistribute.} \\
\bottomrule
\end{tabular}
\end{center}

\noindent
\textbf{Excluded on licence grounds}, recorded so the absence is not mistaken for
oversight: a lung PD-L1 immunohistochemistry set (non-commercial share-alike, and regions
of interest rather than whole slides, with no slide-level tumour proportion score); a
clinico-genomic cohort (non-commercial, no-derivatives); a longitudinal glioma resource
whose agreement forbids commercialising any product containing it; a CNS methylation
classifier whose terms bar revenue-generating use of even its outputs; and one lung
radiomics collection under a non-commercial licence.

\section{External access}

Read-only access can be granted per-prefix without exposing the rest of the bucket, which
also holds the separate lung companion-diagnostic programme. Two options, both tested:

\begin{itemize}[leftmargin=1.4em,nosep]
\item \textbf{Prefix-conditioned IAM} on this bucket grants object reads inside
\texttt{data-request-2026-08/} and denies everything else, including the bucket root.
Verified. Note that object \emph{listing} is a bucket-level permission and cannot be
prefix-scoped, so a conditioned grant yields read-by-path with no browsing --- workable
when paired with this inventory, awkward for exploration.
\item \textbf{A separate share bucket} containing the 27\,GiB analytic layer gives normal
browsing with full isolation, at trivial storage cost. Preferred for an external reviewer.
\end{itemize}

\noindent
Access should be bound to the reviewer's own Google identity rather than granted through
an exported service-account key, so that it is attributable, auditable and revocable per
person. An IAM expiry condition is recommended.

\section*{Bottom line}

\begin{center}
\begin{tabular}{@{}lr@{}}
\toprule
Layer & Size \\
\midrule
A --- DICOM imaging (2{,}031 patients, 17{,}186 series) & 709.8 GiB \\
B --- spatial, multiplex, non-DICOM imaging (5 cohorts) & 540.3 GiB \\
C --- molecular and clinical labels (19 datasets) & $\sim$27 GiB \\
D --- whole-slide imaging + derived features (1{,}186 slides) & 931.4 GiB \\
\midrule
\textbf{Total} & \textbf{$\sim$2{,}209 GiB ($\approx$2.16 TiB)} \\
\bottomrule
\end{tabular}
\end{center}

\noindent
The lung programme is well supported end to end. The glioma programme has a defensible
imaging cohort (630 development, 337 independent validation) and rich molecular reference
tables, but the imaging and the full marker panel belong to \emph{different patients} and
cannot be combined into one cohort --- \S6.3 is the binding constraint. The HCC programme
supports a descriptive lipogenic-phenotype analysis and cannot support a mechanistic one.

\end{document}
